\begin{lemma}
\label{lemma-p-basis}
Let $K/k$ be a field extension. Assume $k$ has characteristic
$p > 0$. Let $\{x_i\}$ be a subset of $K$. The following are equivalent
\begin{enumerate}
\item the elements $\{x_i\}$ are $p$-independent over $k$, and
\item the elements $\text{d}x_i$ are $K$-linearly independent
in $\Omega_{K/k}$.
\end{enumerate}
Any $p$-independent collection can be extended to a $p$-basis of $K$ over $k$.
In particular, the field $K$ has a $p$-basis over $k$.
Moreover, the following are equivalent:
\begin{enumerate}
\item[(a)] $\{x_i\}$ is a $p$-basis of $K$ over $k$, and
\item[(b)] $\text{d}x_i$ is a basis of the $K$-vector space $\Omega_{K/k}$.
\end{enumerate}
\end{lemma}

\begin{proof}
Assume (2) and suppose that $\sum a_E x^E = 0$ is a linear relation
with $a_E \in k K^p$. Let $\theta_i : K \to K$ be a $k$-derivation such that
$\theta_i(x_j) = \delta_{ij}$ (Kronecker delta). Note that any $k$-derivation
of $K$ annihilates $kK^p$. Applying $\theta_i$ to the given relation we
obtain new relations
$$
\sum\nolimits_{E, e_i > 0}
e_i a_E x_1^{e_1}\ldots x_i^{e_i - 1} \ldots x_n^{e_n} = 0
$$
Hence if we pick $\sum a_E x^E$ as the relation with minimal
total degree $|E| = \sum e_i$ for some $a_E \not = 0$, then we
get a contradiction. Hence (1) holds.

\medskip\noindent
If $\{x_i\}$ is a $p$-basis for $K$ over $k$, then
$K \cong kK^p[X_i]/(X_i^p - x_i^p)$. Hence we see that
$\text{d}x_i$ forms a basis for $\Omega_{K/k}$ over $K$.
Thus (a) implies (b).

\medskip\noindent
Let $\{x_i\}$ be a $p$-independent subset of $K$ over $k$. An application
of Zorn's lemma shows that we can enlarge this to a maximal $p$-independent
subset of $K$ over $k$. We claim that any maximal $p$-independent subset
$\{x_i\}$ of $K$ is a $p$-basis of $K$ over $k$. The claim will imply
that (1) implies (2) and establish the existence of $p$-bases.
To prove the claim let $L$ be the subfield of $K$ generated by
$kK^p$ and the $x_i$. We have to show that $L = K$. If $x \in K$
but $x \not \in L$, then $x^p \in L$ and $L(x) \cong L[z]/(z^p - x)$.
Hence $\{x_i\} \cup \{x\}$ is $p$-independent over $k$, a contradiction.

\medskip\noindent
Finally, we have to show that (b) implies (a). By the equivalence of (1)
and (2) we see that $\{x_i\}$ is a maximal $p$-independent subset
of $K$ over $k$. Hence by the claim above it is a $p$-basis.
\end{proof}
